\chapter{绪论}
\label{cha:introduction}

\section{研究背景与意义}
\label{sec:background}

随着移动互联网、智能终端与定位技术的快速发展，基于位置的服务（Location-Based Services, LBS）已经成为人们获取空间信息、发现周边资源和进行城市活动决策的重要入口。地图搜索、导航推荐、本地生活服务、社交签到和旅游出行等应用场景每天都会产生大量带有时空属性的用户行为数据。与传统互联网推荐任务相比，位置相关行为天然包含用户兴趣、时间上下文以及地理空间约束等多重信息，因此如何从这类复杂行为数据中刻画用户移动规律、理解用户即时需求，并进一步实现高质量的位置推荐，已成为推荐系统、数据挖掘和城市计算领域的重要研究课题。

基于位置的社交网络（Location-Based Social Networks, LBSNs）为这一研究提供了典型的数据基础。LBSNs 通常记录用户在真实地理空间中的签到、搜索、点击与访问行为，使得研究者能够从用户轨迹中分析其兴趣偏好、活动区域以及时空迁移模式。围绕 LBSNs 展开的点兴趣（Point-of-Interest, POI）推荐研究，不仅服务于个性化推荐系统本身，也对智慧城市、商业选址、区域活力分析、交通调度和公共服务优化等任务具有现实意义。已有研究表明，LBSNs 与 POI 推荐问题具有信息源异构、数据稀疏、空间依赖显著和上下文敏感性强等特点，因此不能简单地将其视作传统物品推荐的直接延伸，而需要结合时空特征进行针对性建模 \cite{zhao2016survey,sanchez2021survey}。

% 在 POI 推荐任务中，下一个兴趣点预测（Next POI Recommendation）是最具代表性的研究问题之一。该任务旨在根据用户历史访问序列和当前上下文，预测其下一步最可能访问的兴趣点。与一般序列推荐不同，Next POI 推荐不仅要建模用户历史兴趣和行为顺序，还必须考虑 POI 之间的地理距离、区域可达性以及时间规律等因素。现实场景中，用户下一步访问行为往往受到短时活动半径和空间连续性的显著影响，因此单纯依赖语义相似性或协同过滤信号通常难以准确刻画真实移动模式。尤其在城市环境中，不同区域的同类 POI 虽然在类别语义上相近，但其对用户的实际吸引力和可达性可能存在明显差异，这使得地理信息成为 POI 推荐不可忽视的核心因素。

近年来，针对 Next POI 推荐问题，研究者从序列建模、图结构建模、时空联合建模和大模型建模等多个角度进行了探索。早期方法主要依赖马尔可夫链、矩阵分解或基于距离衰减的排序机制，后来逐步发展到循环神经网络、自注意力网络、图神经网络以及超图学习等方法，以提升对用户复杂时空行为的表达能力 \cite{feng2018deepmove,kang2018sasrec,luo2021stan,wang2022graph,yan2023sthgcn}。这些方法在建模用户历史迁移规律方面取得了显著进展，但在 POI 表示学习层面仍然存在明显局限：大量工作仍将 POI 视作原子级离散 ID，或者仅使用类别标签、区域标签等有限侧信息进行编码，导致模型难以充分表达 POI 的深层语义关联及其潜在空间结构。

随着大语言模型（Large Language Models, LLMs）与生成式检索（Generative Retrieval, GR）的兴起，推荐系统研究正逐渐从判别式排序范式转向生成式建模范式。生成式推荐将目标物品预测重构为一个序列生成问题，使模型能够在统一 token 空间中处理用户行为、上下文条件与目标对象之间的关系。相比传统的召回---排序流水线，生成式推荐在统一建模复杂上下文和捕获隐式意图方面展现出较大潜力 \cite{wang2023genrecsurvey,rajput2023tiger,li2023llmgenrecsurvey}。在这一背景下，语义 ID（Semantic ID, SID）逐渐成为生成式推荐中的关键概念。与随机数字 ID 不同，SID 试图为物品构造具有结构性和语义性的离散标识，使相似对象能够在离散 token 空间中共享更长的前缀，从而帮助生成模型学习更稳定的上下文---目标映射关系 \cite{rajput2023tiger}。

% 在 POI 场景中，GNPR-SID 将 Semantic ID 的思想引入 Next POI 推荐任务，提出通过 RQ-VAE 将 POI 的多源特征量化为多层级离散语义标识，再利用大语言模型根据历史访问记录生成下一兴趣点的 SID。该方法为生成式 POI 推荐提供了新的技术路径，也说明了在 POI 任务中构建结构化语义 ID 的可行性 \cite{wang2025gnprsid}。然而，现有基于 SID 的方法仍然面临一个重要问题：POI 本质上同时具有语义属性和地理属性，但现有离散表示学习往往更加关注类别、区域和协同信号，而缺乏对真实地理空间关系的充分建模。若地理信息仅以粗粒度区域标签或简单附加特征的形式参与建模，则最终学得的 SID 更偏向于“语义离散表示”，而难以充分体现 POI 的空间差异和邻近关系。

从应用需求出发，这一问题尤其值得关注。地图搜索和位置推荐场景中的用户查询往往具有较强的空间敏感性和情境依赖性。用户所关注的不只是“该地点属于什么类别”，还包括“它距离当前位置有多远”“是否处于当前活动区域内”“是否适合当前时间和行程安排”等因素。也就是说，对于 Next POI 推荐任务而言，一个高质量的 POI 离散表示不仅应反映类别语义，还应在一定程度上编码地理空间信息，使后续模型在建模用户访问序列时能够更自然地利用空间上下文。

基于上述背景，本文围绕 POI 序列推荐与预测任务展开研究，选择 GNPR-SID 作为基础框架，重点关注如何在 Semantic ID Construction 阶段引入地理信息增强表示，以提升 POI 离散表示的空间感知能力。与直接将 Geohash 或 Geographic ID 显式拼接到大语言模型提示词中不同，本文更关注一种\emph{隐式地理信息注入}思路：通过构造地理信息增强表示，将 POI 的类别语义与地理位置共同编码为连续向量，再输入 RQ-VAE 学习地理增强的 SID。这样得到的 SID 不再只是类别驱动的语义离散编码，而是在表示学习阶段就已经融入了一定的空间结构信息，从而更适合作为生成式 Next POI 推荐的目标表示。

% 从研究意义上看，本文的工作具有三个层面的价值。首先，在理论层面，本文尝试打通“地理信息建模”和“Semantic ID 学习”之间的联系，为生成式 POI 推荐中的语义---空间联合建模提供一种新的思路。其次，在方法层面，本文不通过增加显式输出目标的方式强化地理约束，而是通过优化上游表示学习模块，使地理信息真正影响离散语义表示本身，从而避免输出空间膨胀与生成复杂度上升的问题。最后，在应用层面，本文的研究有助于提升模型对真实地图搜索和位置推荐场景的适应能力，为地图服务、周边推荐和移动出行应用中的个性化地点预测提供方法参考。

% 建议图片1：可以在本节末尾加入“LBSN/Next POI 推荐任务示意图”，
% 展示用户历史访问轨迹、目标时间以及下一 POI 预测的整体场景。
% 例如：用户历史轨迹 + 地图位置分布 + 下一访问点预测箭头。

% % 建议图片2：可以在本节后加入“相关研究发展脉络图”。
% % 例如从 LBSN/POI 推荐 → Next POI 推荐 → 生成式推荐 → Semantic ID 推荐 → 本文方法。
% % 这样能够清楚体现本文工作在整体研究谱系中的位置。


\section{国内外研究现状和相关工作}
\label{sec:related_work}

\subsection{基于位置的社交网络与 POI 推荐研究现状}

基于位置的社交网络（Location-Based Social Networks, LBSNs）为兴趣点推荐研究提供了重要的数据基础。与传统推荐系统不同，LBSNs 中的用户行为不仅包含用户与项目之间的交互关系，还天然伴随着地理坐标、访问时间以及活动区域等上下文信息，因此 POI 推荐问题同时具有推荐系统、时空数据挖掘和城市计算等多重属性。Zhao 等对 LBSNs 场景下的 POI 推荐进行了较早的系统综述，指出该领域的核心难点在于数据稀疏、冷启动、用户偏好动态变化以及地理约束显著等问题 \cite{zhao2016survey}。S{\'a}nchez 等从实验视角进一步总结了 POI 推荐领域的主要研究范式与评价方式，认为 POI 推荐任务不能简单视为传统物品推荐的直接迁移，而必须显式考虑时间与空间因素 \cite{sanchez2021survey}。

围绕这一问题，国内外研究者提出了大量建模方法。早期工作多依赖矩阵分解、概率图模型和基于距离衰减的排序机制，将用户偏好与空间邻近性结合起来进行候选地点打分。这类方法具有一定的可解释性，但面对真实场景下复杂的行为模式时表达能力有限。随着深度学习的发展，研究重心逐渐从浅层统计建模转向深度表示学习。研究者不再仅仅关注“用户过去去过哪些地点”，而是尝试进一步建模“这些地点之间在语义和空间上如何关联”，以及“用户在不同时间与区域中的活动规律如何影响下一次访问决策”。

总体来看，LBSNs 研究的演进为 POI 推荐提供了较为丰富的方法积累，也揭示出该任务的一个核心特点：高质量的地点推荐不仅依赖用户偏好建模，还依赖对 POI 语义结构与地理结构的联合理解。这一点也构成了本文后续研究的基本出发点。

% \subsection{Next POI 推荐研究现状}

% 在 POI 推荐研究中，下一个兴趣点预测（Next POI Recommendation）是最具代表性的任务之一。该任务旨在根据用户历史访问序列预测其下一步最可能访问的地点，因此天然具有序列建模性质。早期相关工作多采用马尔可夫链或时空转移概率建模用户在相邻地点之间的迁移规律，但这类方法通常只能描述较短范围内的局部依赖关系，对于长期兴趣演化和复杂时空上下文的刻画能力较弱。

% 随着深度学习方法的发展，研究者逐步将循环神经网络、自注意力机制和图结构神经网络引入 Next POI 推荐任务。Feng 等提出 DeepMove，通过引入注意力机制和多模态嵌入来建模人类移动轨迹中的长期与短期偏好关系，验证了深度序列模型在移动行为预测任务中的有效性 \cite{feng2018deepmove}。Yang 等提出 Flashback，通过显式回溯历史时空邻域中的相似访问模式来增强下一位置预测能力，使模型在稀疏轨迹场景下仍能够保持较好表现 \cite{yang2020flashback}。Luo 等提出 STAN，在自注意力框架中同时建模时间间隔和空间距离，使模型在注意力权重分配时能够显式感知时空因素，从而提高 Next POI 推荐精度 \cite{luo2021stan}。除这类序列模型外，近年来也有不少工作从图结构视角切入，通过构建用户—POI 图、POI 转移图或时空超图来建模高阶交互关系。例如 Wang 等提出图增强时空网络，对 POI 间的空间关联与顺序依赖进行联合建模 \cite{wang2022graph}；Yan 等则进一步利用时空超图学习增强模型对复杂群组关系和多尺度转移模式的表达能力 \cite{yan2023sthgcn}。

% 尽管这些方法在用户轨迹建模方面取得了显著进展，但大多数研究的重点仍然在于“如何更好地表示用户历史行为序列”，而对“POI 本身如何表示”这一问题关注相对不足。大量方法仍将 POI 视作离散原子 ID，或仅使用简单的类别标签和区域标签进行辅助建模。这意味着即使模型能够较好地学习用户序列模式，也未必真正理解不同 POI 在语义与空间上的细粒度差异。对于 Next POI 推荐任务而言，这种表示层面的不足会直接限制模型性能的进一步提升。

\subsection{序列推荐与大模型推荐研究现状}

从更一般的推荐系统研究来看，Next POI 推荐可以被看作序列推荐在位置服务场景中的特殊实例。序列推荐的核心目标是根据用户历史行为序列预测未来偏好，其关键在于如何建模用户兴趣随时间的动态演化。Kang 等提出 SASRec，首次系统展示了自注意力机制在序列推荐任务中的优势，相比传统循环神经网络模型，自注意力结构能够更灵活地捕捉长距离依赖关系 \cite{kang2018sasrec}。此后，大量研究进一步围绕序列推荐展开，在自监督学习、多行为建模、跨域推荐和鲁棒表示学习等方面取得了丰富成果。

近年来，大语言模型的发展又推动了推荐系统研究范式的进一步转变。Wang 等在生成式推荐综述中指出，推荐系统正在从以“打分排序”为核心的判别式框架，逐步走向以“条件生成目标对象”为核心的生成式框架 \cite{wang2023genrecsurvey}。Li 等进一步总结了大语言模型在生成式推荐中的潜力，认为 LLM 在处理复杂上下文、统一建模用户行为和提高推荐可解释性方面展现出显著优势 \cite{li2023llmgenrecsurvey}。对于 POI 场景而言，历史访问记录、查询文本、时间上下文和地点表示都可以自然组织成结构化序列，因此大语言模型为 Next POI 推荐任务提供了新的建模可能。

但是，大语言模型虽然具备较强的上下文建模能力，其最终效果仍然高度依赖目标对象的表示形式。如果目标仍然只是无意义的原子 ID，那么即使模型结构再强，也很难学习到稳定而有泛化能力的上下文—目标映射关系。因此，在大模型推荐研究逐渐兴起的背景下，如何为推荐对象设计更适合生成模型学习的离散表示，成为了一个重要研究问题。

\subsection{Semantic ID 推荐研究现状}

为了缓解生成式推荐中随机原子 ID 难以学习的问题，Semantic ID 逐渐成为推荐系统和生成式检索中的重要研究方向。Semantic ID 的基本思想是通过量化编码、层次聚类或残差向量量化等方式，将对象的连续表示映射为多层级离散 token 序列，使得语义相似的对象在离散空间中共享更长的前缀，从而帮助生成模型学习对象之间的潜在结构关系。

Rajput 等在 TIGER 中利用残差量化方法为推荐对象构造语义 ID，展示了这种结构化离散表示在生成式推荐中的有效性 \cite{rajput2023tiger}。Hou 等进一步研究了更长语义 ID 的并行生成问题，说明随着生成任务复杂度提升，如何构造更具可学习性的语义标识已经成为生成式推荐中的重要议题 \cite{hou2025longsid}。在 POI 推荐领域，Wang 等提出 GNPR-SID，将 Semantic ID 引入 Next POI 推荐任务，通过整合类别信息、区域信息、时间信号与协同访问特征构造连续 POI 表示，并利用 RQ-VAE 学习多层级 SID，再通过大语言模型完成下一兴趣点的生成式预测 \cite{wang2025gnprsid}。该方法表明，相比原始数字 ID，具有语义结构的离散表示更适合用于 POI 序列推荐任务。

然而，现有 Semantic ID 方法普遍存在一个共同局限，即它们更侧重于对象的语义属性和协同行为模式，而对对象所携带的其他结构信息建模不足。对于一般推荐场景而言，这种语义导向的设计已经能够带来明显提升；但对于 POI 推荐任务来说，地点不仅具有类别和文本语义，还具有强烈的空间属性。如果 SID 的学习主要由类别、区域标签和协同信号驱动，那么最终学得的离散标识很可能更接近“语义相似对象的离散聚类结果”，而不能充分反映空间上相近或空间上差异显著的 POI 之间的关系。因此，在 POI 场景中，如何构造兼具语义结构与空间感知能力的 Semantic ID，仍然是一个值得进一步研究的问题。

\subsection{地理信息建模研究现状}

地理信息建模是 POI 推荐区别于一般推荐任务的关键所在。现有研究中，地理信息的引入方式大体可以分为三类。第一类方法将经纬度、地理距离或空间邻近性作为附加数值特征直接输入模型。这类方法实现简单，能够一定程度上约束推荐结果的空间合理性，但往往难以与语义表示形成统一表示空间。第二类方法将地理信息作为时空联合建模的一部分，在注意力权重、图结构边关系或轨迹转移建模中显式引入空间距离，以增强模型对用户移动行为的表达能力。第三类方法则进一步尝试构造地理嵌入、区域图或地理编码，使空间信息不再只是被动附加，而成为 POI 表示的一部分。

在 POI 检索与地图搜索场景中，Zhao 等提出 PALM，将语义相似性与地理相关性联合用于 query-POI 相关性建模，说明地理因素对位置场景中的排序性能具有重要影响 \cite{zhao2019palm}。Huang 等提出 HGAMN，通过异构图注意力网络对查询、POI 及多种上下文关系进行联合建模，在地图检索场景中取得了良好效果 \cite{huang2021hgamn}。在更偏生成式的地图搜索研究中，GenPOI 进一步提出将地理显式标识与地理位置嵌入共同融入 POI tokenization 过程，试图让大语言模型同时理解“POI 是什么”和“POI 在哪里”:contentReference[oaicite:2]{index=2}。其中，GeoPE 通过参考点方位角与嵌入旋转方式隐式注入空间信息，说明地理信息不仅可以显式编码为 geohash，也可以通过连续表示变换方式融入潜在语义空间:contentReference[oaicite:3]{index=3}。

总体来看，现有工作已经充分证明地理信息在 POI 推荐和 POI 检索任务中的重要性，但多数方法仍然是在推荐预测阶段或检索排序阶段引入地理约束，而较少真正作用于上游的离散语义表示学习过程。也就是说，现有研究普遍回答了“如何让模型在预测时使用地理信息”，却较少回答“如何让地理信息直接参与 Semantic ID 的形成”。这一空缺恰恰构成了本文工作的研究切入点。

\section{本文研究内容与章节安排}
\label{sec:main_work}

基于上述研究背景与现状分析可以看出，Next POI 推荐任务已经从传统的时空建模逐步发展到结合大语言模型和生成式推荐的新阶段，但 POI 表示学习仍然是制约模型性能的关键因素之一。特别是在 GNPR-SID 这类基于 Semantic ID 的方法中，SID 的质量直接决定了后续生成模型的学习难度和预测效果。如果 SID 仅由类别、区域和协同信号学习得到，而缺乏对真实地理空间关系的充分建模，那么模型很难在输出阶段真正理解同类 POI 的空间差异和用户活动区域约束。

围绕这一问题，本文在 GNPR-SID 框架基础上开展研究，重点关注如何在 Semantic ID Construction 阶段引入地理信息增强表示，使地理信息能够在 POI 离散表示学习过程中发挥作用。本文首先对原始签到数据进行整理与清洗，构建用户访问序列和 POI 属性表；随后，围绕 POI 的类别名称和地理坐标构造地理信息增强表示，并利用 RQ-VAE 学习地理增强的 SID；在此基础上，进一步沿用 GNPR-SID 的生成式推荐框架，将学习得到的 SID 组织为结构化 token 序列，输入大语言模型进行监督微调，以实现 Next POI 推荐。

需要说明的是，本文在方法探索过程中也尝试过将 Geohash 等地理编码显式加入生成式推荐模块，即直接把地理信息注入提示词或输出目标中，希望通过显式 token 生成强化模型对空间信息的理解。然而实验表明，这种方式虽然能够让大语言模型“看到”地理信息，却也显著增加了输入输出长度和生成复杂度，导致推理开销变大，且对核心 SID 预测任务的帮助有限。相比之下，在表示学习阶段通过 GeoPE 等机制将地理信息隐式注入连续表示，再由 RQ-VAE 学习出地理增强 SID，往往能够在预测效果与计算代价之间取得更好的平衡。因此，本文最终选择以\emph{隐式地理信息注入}作为核心方法路线。

全文的组织结构如下。第二章介绍本文涉及的相关理论与关键技术，包括 POI 序列预测、RQ-VAE、大语言模型与参数高效微调以及地理空间散列编码等内容。第三章在 GNPR-SID 原始框架的基础上，详细介绍本文提出的基于地理信息增强表示的改进方法，重点说明显式与隐式地理注入方案的差异，以及 GeoPE 与地理增强 SID 的学习过程。第四章给出实验设计与结果分析，从总体性能、消融实验和方法讨论等角度验证本文方法的有效性与局限性。最后，第五章对全文工作进行总结，并对后续可能的改进方向进行展望。